\usepackage[utf8]{inputenc}
\usepackage{multicol}
\usepackage{xcolor}
\usepackage{subfigure}
\usepackage[spanish,es-tabla]{babel}
\usepackage[utf8]{inputenc}
\usepackage{graphicx}
\usepackage{titlesec}
\usepackage[bookmarks,breaklinks,colorlinks=true,allcolors=blue]{hyperref}
\usepackage{listings}
\usepackage{inconsolata}
\usepackage{float}
% \usepackage{mathpazo} % Fuente Palatino -- comentado SOLO para preview local (falta el paquete "palatino" en tu BasicTeX). Descomentar al usar Overleaf, o tras "sudo tlmgr install palatino".
\usepackage[labelfont=bf]{caption}

\usepackage[square,numbers]{natbib}
\usepackage[nottoc,notlof,notlot]{tocbibind}  % Mete la bibliografía como capítulo en la TOC, los parámetros excluyen los otros índices de aparecer también
\usepackage{geometry}
\usepackage{amsmath}
\usepackage{amssymb} % \checkmark, entre otros símbolos usados en tablas
\usepackage{booktabs} % Estilo de tablas con \toprule/\midrule/\bottomrule
\usepackage{colortbl} % \rowcolor, \cellcolor -- cabeceras de tabla resaltadas
\usepackage{multirow} % Celdas que ocupan varias filas en tablas (p.ej. agrupar hiperparámetros por algoritmo)
\usepackage{longtable} % Tablas de cuadrícula completa que pueden partirse entre páginas
\usepackage{array} % Modificadores de columna (p.ej. centrado en columnas p{}) en tablas
\usepackage{tikz} % Diagramas EDT y de arquitectura dibujados directamente en LaTeX
\usetikzlibrary{arrows.meta} % Puntas de flecha modernas (p.ej. -{Latex[...]}) en los diagramas tikz
\usetikzlibrary{calc} % Coordenadas calculadas (p.ej. punto medio entre dos nodos) en los diagramas tikz
\usetikzlibrary{shapes.geometric} % Elipses (nodos de caso de uso UML) en los diagramas tikz
\usepackage[section]{placeins} % Fuerza a colocar todas las tablas/figuras pendientes antes de cada nueva sección, evitando que "adelanten" al texto siguiente
\usepackage{parskip}
\usepackage[official]{eurosym}
\usepackage{todonotes}
\usepackage{csquotes}
\usepackage{tocbasic}  % Estilos de la TOC
